\documentclass[10pt,twocolumn,letterpaper]{article}

\usepackage{iccv}
\usepackage{times}
\usepackage{epsfig}
\usepackage{graphicx}
\usepackage{amsmath}
\usepackage{amssymb}
\usepackage{subcaption}
\usepackage{multirow}
% Include other packages here, before hyperref.

% If you comment hyperref and then uncomment it, you should delete
% egpaper.aux before re-running latex.  (Or just hit 'q' on the first latex
% run, let it finish, and you should be clear).
\usepackage[pagebackref=true,breaklinks=true,letterpaper=true,colorlinks,bookmarks=false]{hyperref}

% \iccvfinalcopy % *** Uncomment this line for the final submission

% \def\iccvPaperID{9028} % *** Enter the ICCV Paper ID here
\def\httilde{\mbox{\tt\raisebox{-.5ex}{\symbol{126}}}}

% Pages are numbered in submission mode, and unnumbered in camera-ready
\iccvfinalcopy
\ificcvfinal\pagestyle{empty}\fi

\newcommand{\prashant}[1]{{\color{red} \textbf{(Prashant: #1})}}

\newcommand{\dslr}{\texttt{\textbf{MOVESe}}}
\newcommand{\dslruda}{\texttt{\textbf{MOV-UDA}}}
\newcommand{\CacciaAE}{\texttt{\textbf{Caccia-AE}}}
\newcommand{\CacciaVAE}{\texttt{\textbf{Caccia-VAE}}} 
\newcommand{\CacciaGAN}{\texttt{\textbf{Caccia-GAN}}}

\newcommand{\Coarse}{\texttt{\textbf{Coarse-Net}}}

% \newcommand{DSLR}{\texttt{\textbf{DSLR}}}
\newcommand{\MYhref}[3][black]{\href{#2}{\color{#1}{#3}}}

\newcommand{\siv}{\texttt{\textbf{ShapeInversion}}}

\begin{document}

%%%%%%%%% TITLE
\title{MOVESe: \underline{MOV}abl\underline{E} and Moving LiDAR Scene \underline{Se}gmentation with Improved Navigation in Seg-label free settings}

\author{Prashant Kumar\textsuperscript{\rm \footnotemark[1]1}, Onkar Susladkar\textsuperscript{\rm \footnotemark[1]2}, Dhruv Makwana\textsuperscript{\rm \footnotemark[1]2}, Anurag Mittal\textsuperscript{\rm \footnotemark[1]1}, Prem Kumar Kalra\textsuperscript{\rm \footnotemark[1]1}\\
\textsuperscript{\rm \footnotemark[1]1}Indian Institute of Technology, Delhi, \textsuperscript{\rm \footnotemark[1]2}Independent Researcher \\
}

\maketitle
% Remove page # from the first page of camera-ready.
\ificcvfinal\thispagestyle{empty}\fi

%%%%%%%%% ABSTRACT
\begin{abstract}
Accurate detection of movable and moving objects in LiDAR is of vital importance for navigation. Most existing works focus on extracting and removing moving objects during navigation. Movable objects like pedestrians, parked vehicles, etc. although static may move in the future. This leads to erroneous navigation and accidents. In such cases, it becomes necessary to detect potentially movable objects. To this end, we present a learning-based approach that segments movable and moving objects by generating static parts of scenes that are otherwise occluded. Our model performs superior to existing baselines on static LiDAR reconstructions using 3 datasets including a challenging sparse industrial dataset. We achieve this without the assistance of any segmentation labels because such labels might not always be available for less popular yet important settings like industrial environments. The non-movable static parts of the scene generated by our model are of vital importance for downstream navigation for SLAM. The movable objects detected by our model can be fed to a downstream 3D detector for aiding navigation. 
Though we do not use segmentation, we evaluate our method against navigation baselines that use it to remove dynamic objects for SLAM. 
Through extensive experiments on several datasets, we showcase that our model surpasses these baselines on navigation.
\end{abstract}

%%%%%%%%% BODY TEXT
\section{Introduction}
\label{introduction}
\begin{figure}[htbp]
\centering
\includegraphics[width=1.0\linewidth, height=0.9\linewidth]{first-page.pdf}  
  \caption{Top view of 3 contiguous interleaving scans (left to right) with movable objects (pink box) and moving objects (blue box). \textbf{Top}: Dynamic LiDAR scans \textbf{Middle}: Reconstructed static scans with all dynamic objects removed. \textbf{Bottom}: Top and middle scans used to generate moving and movable objects (Black dot indicates origin (0,0)).  Note the changing positions of moving objects while the movable object remains stationary.  \textit{In the 3rd column, object 3 changes from moving to movable.} Video demo in Supp.}
\label{fig:movable-first}
\end{figure}

Handling moving and movable objects is an important problem that has major applications in navigation and allied tasks. Most existing techniques focus on extracting and removing moving objects while using only static structures for navigation \cite{ruchti2018mapping,chen2021moving,nishio2022moving,pfreundschuh2021dynamic,tian2022dl}. They do so by using labeled segmentation and object information to train a model to identify moving objects. However, everyday navigation scenarios come across \textit{potentially movable} objects. Such objects have the capacity to move but may appear stationary at a given instant and hence are a potential hazard to downstream navigation-based tasks. When moving, these objects also add to the moving parts of a scene, thereby adding to the Absolute Translation Error \cite{sturm2012benchmark} and Relative Pose Error \cite{sturm2012benchmark} during pose estimation  \cite{kumar2021dslr}. Vaquero et. al. \cite{vaquero2019improving}  propose to remove all moving and movable objects for SLAM, with the help of KITTI segmentation labels \cite{behley2019semantickitti}. However, such dataset are difficult to collect and label for less popular for important settings like industrial environments.

Central to our task of discriminating \textit{dynamic objects (moving and movable)} from \textit{static (non-movable)} objects is the design of the discriminator. For this task, we require the availability of LiDAR scans that are only composed of only static structures in a given environment, which we refer to as \textit{static scans}. This is important in order for our model to be able to distinguish dynamic objects from immovable objects. Using a standard discriminator that tells apart completely static scans from dynamism-induced scans (\textit{dynamic scans}), suffers from a \textit{simplicity bias} \cite{shah2020pitfalls}. The discriminator needs to only pay attention to the \textit{simplest features} to use for discrimination \textit{i.e.(moving objects)}, to distinguish them from static scans. It can prematurely saturate by learning on to these simplest features at the
earliest. Hence, movable objects do not contribute significantly to the learning. as the model may never need to use the more complex movable object features (as they are difficult to discover). This leads to  suboptimal static translations and inaccurate segmentation (Fig. \ref{fig:ati-moving-movable}).

In this work, we address the aforementioned issues and propose the following contributions:
\begin{itemize}

    \item We present \dslr{}, a GAN-based adversarial model for \textit{generating static background} of occluded parts of a dynamic LiDAR scan and henceforth \textit{segmenting out moving and movable  parts of scenes}. This is achieved in a segmentation label-free manner. \dslr{} uses a novel couple discriminator to discriminate carefully chosen pairs of static and dynamic scans. The novel design choice of the input helps our model to explicitly focus its attention on moving as well as movable objects, without any external supervision.
 
    \item To reinforce the attention of our couple discriminator on moving and movable objects, we design a novel contrastive strategy. Here, we utilize paired correspondence information between static and dynamic scans to design hard negative triplets. This assists our discriminator and allows it to pay explicit attention to moving and movable objects.

    \item We quantitatively illustrate the usefulness of static reconstructions for navigation using SLAM. We also qualitatively demonstrate its usefulness in segmenting out movable as well as moving objects from dynamic LiDAR scans. Our strategy of using the generated static structures for SLAM performs superior to existing approaches on $3$ varied datasets and comfortably surpasses several important baselines. Surprisingly, \textit{our model also performs better than models that use ground truth segmentation labels information} to remove dynamic objects/points, thereby indicating the importance of reconstructed static structures for SLAM.
    
   \item To handle datasets that lack the static-dynamic pair mapping information, we propose \dslruda{}, that uses Unsupervised Domain adaptation. \dslruda{} also demonstrates superior performance on SLAM.     
\end{itemize}

\section{Related Work}
\label{related}
 Generative modeling for LiDAR was first explored by Caccia et. al. \cite{caccia2018deep} where they use LiDAR used range image-based representation for  LiDAR reconstruction and generation. Generating static occluded parts of scenes has been explored for several modalities: images, point clouds, and recently LiDAR scans. For images, Bescos et. al. \cite{bescos2019empty} erase dynamic objects, with the help of a semantic segmentation network, from images obtained from CARLA simulator \cite{dosovitskiy2017carla}, and on Cityscapes via synthetic-to-real transfer. Wang et al. \cite{wang2022coarse}, propose a coarse-to-fine network for dynamic-to-static image translation, \textit{without any segmentation information} and show impressive results.  With respect to point clouds Chen et. al. \cite{chen2019unpaired} attempt to complete partial point cloud scans of objects \cite{chang2017matterport3d,dai2017scannet} to their original shape using generative models. Further, Wu et. al. \cite{wu2020multimodal} provide multi-modal reconstructed shapes for a given incomplete point cloud using conditional generative models. Another closely related work by Zhang et. al. \cite{zhang2021unsupervised}, uses GAN inversion for point cloud completion and shows remarkable results on standard shape datasets, and we use it as a baseline. Another work by Kumar et al. \cite{kumar2021dslr}, achieves static translations for LiDAR scans using an adversarially trained model. Although they implicitly don't use segmentation, they have 2 variants that use segmentation information for static translation as well as for downstream SLAM-based navigation. Further Yang et. al. \cite{yang2021semantic}, Cheng et al. \cite{cheng2020s3cnet} also approach segmentation-assited scene completion for LiDAR.
 

Several label-free static translation methods as discussed above, focus on enhancing the power of the generator. Junzhe et al. \cite{zhang2021unsupervised} use a novel loss, PatchVariance, and a novel masking technique, k-mask for the generator. Similarly, Teng et al. \cite{wang2022coarse}, uses an enhanced coarse-to-fine generator for their goal. On the other hand, our focus is to enable the discriminator to learn discriminatory features using moving as well as movable objects.  \\
Vaquero et. al. \cite{vaquero2019improving} use KITTI semantic segmentation information for segmenting all movable and moving objects for LiDAR. To the best of our knowledge, there do not exist any segmentation label-free methods for movable object segmentation. \\ 

\textbf{Contrastive Learning:}
Contrastive Learning learns a rich latent representation of data via grouping samples with the same labels while distancing out dissimilar labeled samples.
Some well-known contrastive losses include Maximum Margin Contrastive Loss \cite{hadsell2006dimensionality}, Triplet Loss \cite{weinberger2009distance} and multi-class N-pair loss \cite{sohn2016improved}.
Triplet Loss \cite{weinberger2009distance} operates on a triplet of latent vectors. Two vectors (the anchor and positive) share the same label while the third (negative) belongs to a different label. The objective of triplet loss is to keep the positive sample within a \textit{m-neighborhood of the anchor while forcing the negative sample} out of the m-neighborhood. 
The multi-class N-pair loss  is a generalization of the triplet loss.
Contrastive Learning with hard negatives \cite{robinson2020contrastive} involves selecting negative samples that are difficult to distinguish from the anchor sample. Several approaches \cite{cao2020parametric,chuang2020debiased,iscen2018mining,wu2020mutual,xie2022delving,ho2020contrastive} deploy hard negative mining for contrastive learning in order to learn better representations for downstream tasks.

\section{Problem Description}
We refer to a scan pair as a \textit{corresponding static-dynamic LiDAR} scan pair where both the scans share \textit{exactly the same surrounding, except for dynamic occlusions}. Assuming we have a set of static scans $St$, and corresponding dynamic scans denoted as $Dy$.
We learn to map a point ($r_d$) on the dynamic LiDAR latent manifold to its corresponding point on the static LiDAR latent manifold ($r_s$) such that $r_s$ gives the corresponding static LiDAR by generating static background of the occluded regions.\\ 
  We resort to range image-based LiDAR preprocessing \cite{caccia2018deep}. We convert LiDAR scans to range images, which allows us to use image-based processing on LiDAR efficiently.

\subsubsection{\dslr{} Formulation:}
For our static translation, we consider a conditional GAN architecture. We use the Encoder ($E_\phi$) and Decoder ($D_\theta$) from \cite{caccia2018deep}. Together these form the generator (\textit{G}) for our model.\\ Given $d_i\in Dy$ is input LiDAR scan, the goal of $G$ is to generate $s_i$ ($s_i \in St$) (${s_i}$ is the corresponding static of $d_i$). 
\\



\textbf{Couple Discriminator(DI):} A  challenge in discriminating static from dynamic frames is that both are generated from the same distribution. Dynamic LiDAR scans are essentially static LiDAR with noise in the form of dynamic objects. The latent feature for both types of scans consists of attributes that are mostly, \textit{homogeneous and indistinguishable} because both share similar LiDAR scenes and structure. Further, attributes that discriminate them are scarce because dynamic occlusions (due to moving and movable objects) occupy a small percentage of points in a LiDAR scan. Therefore, there exists an imbalance w.r.t. the discriminatory features. \\
\textit{We use a novel couple discriminator} (\textit{DI}) for our proposed GAN, which operates on the latent space of the LiDAR scans. Let ($r_{s_i}$, $r_{s_j}$, $r_{d_k}$) be latent representation for the triplet (${s_i}$, ${s_j}$, ${d_k}$).
DI takes as input a pair of LiDAR scan latent representations. ($r_{s_j}$, $r_{s_i}$) is a labelled as real while ($r_{d_k}$, $r_{s_i}$) as fake. 

Such a setting enables DI to focus on discriminatory features between static and dynamic scans. In order to classify the above pairs as real and fake respectively, the \textit{discriminator is forced to amplify discriminatory features in dynamic frames}. Instead of competing static representations against dynamic (as in a standard discriminator), here discrimination occurs between \textit{similarities in static scans v/s differences between static and dynamic scans}. Thus, DI can explicitly focus on  dynamic occlusions(due to moving and movable objects) in dynamic scans. The above behavior is assisted by the fact that the LiDAR scans in \textit{St} are completely devoid of moving and movable objects, which are present in \textit{Dy}.
 This factor helps us in segmenting out movable as well as moving objects using generated static LiDAR scans. 
and also on the static scenes behind those occlusions in static scans.  It also helps DI \textit{differentiate between occlusions due to static and dynamic objects}, which helps G \textit{remove dynamic objects and dynamic occlusion while sparing static ones}). 

Further to stabilize the training of DI, we replace the negative log-likelihood objective with the least square objective \cite{mao2017least}. The idea behind using Least Square GANs is 2 fold: first is to ensure that DI receives smooth and non-saturating gradients, thereby allowing \textit{G} which is also trained via gradients of DI to be trained effectively. Secondly, given gradients could be weak due to the structural similarity between static and dynamic frames, we'd even want to \textit{penalize LiDAR pairs that are far from the discriminator decision boundary,  though on the correct side}. This will allow $G$ to generate static reconstruction close to the decision boundary, which in the ideal case should go across the manifold of real data to avoid learning saturation. \cite{mao2017least}.  This helps generate more gradients while updating $G$, thus allowing translated samples to be closer to the real manifold of static scans.\\
\textbf{Contrastive Learning on Couple Discriminator:} Dynamic and static scans appear to be generated from similar distributions and share similar features (except for moving and movable objects). Thus, learning discriminatory features between them is non-trivial. To aid DI in paying attention only to differences between them, we use contrastive loss to \textit{contrast dynamic scans from static} ones while bringing static scans closer. Here we use a careful choice of static and dynamic scans using \textit{hard negative mining}. Consider the triplet above $(s_i\in St, s_j\in St, d_k\in Dy)$ where we now constrain \textit{$d_k$ to be the corresponding static scan for $s_{i}$} (Thus $d_{k}$ is now $d_{i}$). We form couples $(s_j, s_i)$ and $(d_i, s_i)$ for the discriminator loss.\\


Consider $s_i$ as the anchor, $s_j$ as positive, and $d_i$ \textit{(corresponding dynamic/hard negative for $s_i$)} as the negative scan. Thus, the \textit{negative scan is more closer to the anchor than the positive} in the latent space.
Triplet loss brings static LiDAR couple $(s_j, s_i)$ close together within a \textit{m-neighborhood} of the anchor scan ($s_j$) in a metric space while distancing out negative and anchor LiDAR couple $(d_i, s_i)$ i.e. pushing negative scan ($d_i$) outside the m-neighborhood, m being a hyper-parameter known as margin.\\ We now explain the reasons behind choosing $d_i$ as the hard negative for the anchor $s_i$ in the triplet. Given $s_i$ and $d_i$ are perceived in the same environment and location, they share the same static structures. The only difference between them (moving and movable objects) accounts for the discriminatory features. Contrasting them \textit{w.r.t} each other forces DI to only focus on the moving and movable objects (\textit{the only available discriminatory features}) in $d_i$, and on the corresponding static background features in static scans($s_i$). This ensures that the signals that flow to $G$ explicitly focus on the above structures, thereby helping identify moving and movable objects.
This also prevents \textit{occlusion due to static objects from being mistaken as dynamic}, thereby overcoming limitations of existing methods (Section \ref{related}). These improvements strengthen discriminatory signals for DI while also helping $G$ focus on only dynamic occlusion and objects explicitly for reconstructing background static scenes. Given \textit{d} as a metric, contrastive triplet loss is:
\begin{equation}
        =max(0, d( s_i, s_j)  - d(d_i, s_i) + m)
\end{equation}


\subsection{Training Setup for Conditional Translation}
\begin{figure}[htbp]
  \centering
  \includegraphics[width=1.0\linewidth, height=0.9\linewidth]{ECCV-Final.pdf}  
  \caption{Training Setup for our model.}
\label{fig:carlagan}
\end{figure}

We detail the training setup for our model (Figure \ref{fig:carlagan}) that helps to achieve the dynamic to static translation and thereby segmentation of moving and movable objects.\\
 \textit{Discriminator}: We use a MLP-based sequential neural network consisting of 6 layers with sigmoid non-linearity. As detailed above, for training, we consider a triplet of the form $(s_i\, s_j, d_i)$ (shown in Fig. \ref{fig:carlagan}) for training DI. We generate latent vectors of the triplet using the encoder ($E_{\phi}$ of G), and create pairs($r_{s_j}$, $r_{s_i}$) and ($r_{d_i}$, $r_{s_i}$). The former is labeled as real while the latter as fake by our discriminator. \\

\textit{Generator:} The generator is fed with a dynamic scan ($d_i$) and is reconstructed using the corresponding static frame, $s_i$ (Figure \ref{fig:carlagan}). Also, note that the latent representation ($r_{d_i}$) of $G$ is a constituent of the pair which is labeled as fake by DI.
Thus, the generator is trained using the signals received from DI, as well the backprop signals from the decoder($D_{\theta}$).
\textit{For a detailed description of the Loss function for the G and DI, please refer to Supp}.

Now we detail the training process of our GAN.
Given a batch of triplets $(s_i,s_j,d_i)$, we use Encoder ($E_{\phi}$) of G, to extract latent representations for the full batch of triplets, i.e. $(r_{s_i},r_{s_j},r_{d_i})$. 
Using these we create the real and fake pairs, ($r_{s_j}$, $r_{s_i}$) and ($r_{d_i}$, $r_{s_i}$) respectively.\\
DI is trained to classify these pairs correctly.
However, given a dynamic frame ($d_i$) as input, the generator $G$ tries to produce a latent representation ($r_{d_i}$), such that the pair ($r_{d_i}$, $r_{s_i}$) is labeled as real instead of fake by the discriminator. In order to fool the discriminator, \textit{G} must ensure that it generates latent representations on the static manifold (i.e similar to static LiDAR scans) while it is given a dynamic frame as input. This strategy helps us to achieve our translation.

\subsection{\dslruda{} }

One of the current challenges for our approach is the unavailability of corresponding ground truth static for dynamic LiDAR scans in some real-world scenarios. Although ARD-16, a real-world industrial dataset, has correspondence information available, KITTI dataset lacks it. Therefore, to adapt models trained on paired datasets to unpaired settings, we consider using Unsupervised Domain Adaptation. We follow Borgwardt et. al.\cite{Borgwardt2006IntegratingSB} and integrate the domain discrepancy loss between CARLA-64 and KITTI-64 in the latent space. This projects data points from both domains in reproducible Kernel Hilbert space followed by computing the average difference between them. For adapting CARLA domain scans, to the KITTI domain (Refer Figure \ref{fig:ARCHDIAG}), we consider the same architecture for the GAN as discussed above. The only additions are :-  \\$\xrightarrow[]{}$ KITTI dynamic scan latent representations ($E(k_l)$ or $r_{k_l}$) is paired with $r_{s_j}$. This pair ($r_{k_l}, r_{s_j}$) is fed to $DI$, and is labelled as fake. \\
$\xrightarrow[]{}$ Domain loss is calculated between $r_{k_l}$ (KITTI scan) and $r_{s_i}$ (static LiDAR scan) and added to generator loss.\\ 
$\xrightarrow[]{}$ KITTI dynamic scan is reconstructed at the KITTI generator $G(\gamma, \lambda)$. This helps the KITTI scan to retain static structures. \textit{For a detailed description on the Loss function, please refer Supp}.

\section{Experiments}
We evaluate our proposed model for static translations on 3 different datasets including an \textit{indoor and outdoor dataset} using 5 baselines that show strong modeling performance on LiDAR for static translation and have been adapted for our task. We show quantitative results for static translations (\ref{sec:chamfer-nu}) and the benefit of predicted static structures in SLAM (\ref{sec:slamindyenv}). We further show qualitative results for moving and non-moving object segmentation for LiDAR (\ref{sec:segm_mova_movi}).
\subsection{Datasets:}

\textbf{ARD-16}: ARD-16 \cite{kumar2021dslr,dslr-git} is a real-world industrial  LiDAR dataset collected using a 16-beam VLP-Puck LiDAR sensor \cite{velo16}. Like CARLA-64, it also has paired correspondence information available. This dataset is challenging given it is 4$\times$ sparse compared to the other datasets and has very few points falling on static structures as well as on dynamic objects. This makes it difficult to learn occlusion and static structure-based features. We use this dataset to show that our model works reasonably well on challenging sparse datasets in the absence of segmentation information. For training and testing \dslr{} we follow the same protocol as Kumar et. al. \cite{kumar2021dslr}. For SLAM, we use the single run released \cite{dslr-git}, which is longer and combines all sequences and scenes available in previous testing runs.\\
\textbf{KITTI}: We use KITTI dataset to experiment with our model using Unsupervised Domain Adaptation on a real-world 64-beam dataset. For KITTI, less than 7\% of the available scans are static scans (no dynamic occlusion) and without any correspondence information available. Therefore, we are not able to use standard metrics like Chamfer's Distance for evaluation on KITTI. Thus, we resort to LiDAR Quality Index (LQI) which has been shown to positively correlate with Chamfer's Distance and has been used in such scenarios \cite{kumar2021dslr}. More details in Supp. 


For SLAM, we use the KITTI Odometry dataset sequences from 00-10 which have ground truth poses available for SLAM. We show the results on all runs to study the behavior of our model.\\
 \textbf{CARLA-64} : CARLA-64 is extensive, simulated dataset collected using CARLA simulator\cite{dosovitskiy2017carla}. It is collected using settings of VLP-64 Velodyne LiDAR on a virtual city map\cite{kumar2021dslr}. We perform our experiments on the latest version of the dataset and follow the train/test protocol  as mentioned here\cite{dslr-git}. For SLAM, we use the 4 available dynamic runs for CARLA-64.

 One of the motivations behind using CARLA is the \textit{availability of paired correspondence} between static and dynamic scans which is vital for such a framework to succeed. Such paired dataset is very scarce in literature. However, \cite{kumar2021dslr} show that it is viable and easy to collect such dataset in urban and industrial settings.\\
\begin{figure}
  \centering
  \includegraphics[width=1.0\linewidth, height=0.8\linewidth]{ECCV-Kitti-Final.pdf}  
  \caption{\dslruda{} training using MMD Loss for KITTI.}
\label{fig:ARCHDIAG}
\end{figure}



\subsection{Baselines}
For our baselines, we choose models that have shown impressive modeling performance on LiDAR datasets for static translations, without using segmentation.  From the available literature, we conclude that learning-based methods \cite{kendall2015posenet,krull2015learning,bescos2019empty} outperform non-learning based methods.

The following methods show reasonable to strong performance on LiDAR for our task. Thus, we use them as our baselines: 1.\CacciaAE{} (C-AE) 2. \CacciaGAN{} (C-GAN) \cite {caccia2018deep} 3.\Coarse (C-NET) \cite{wang2022coarse} 4. \siv{}(S-INV) \cite{zhang2021unsupervised} 5. DSLR \cite{kumar2021dslr}. For more details on baselines, please refer Supp.
\subsection{Evaluation Metrics:}
\label{metrics} For evaluation, we use Chamfer Distance \cite {fan2017point} and Earth Mover's Distance \cite {emd-cd} for evaluating the closeness of the translated dynamic scans to corresponding static LiDAR scans for CARLA-64 and ARD-16  datasets. For the KITTI dataset, we adopt LQI metric \cite{kumar2021dslr}. For details, please refer supp.


\begin{figure}
     \centering
     \begin{subfigure}[b]{0.30\linewidth}
         \centering
         \includegraphics[width=\textwidth, height=0.5\linewidth]{original_955_colored_updated.png}
         \label{fig:sub-tenth}
     \end{subfigure}
     \begin{subfigure}[b]{0.30\linewidth}
        \centering
        \includegraphics[width=\textwidth, height=0.5\linewidth]{ours_955_colored_updated.png}
        \label{fig:sub-twelvth}
    \end{subfigure}
    \begin{subfigure}[b]{0.30\linewidth}
        \centering
        \includegraphics[width=\textwidth, height=0.5\linewidth]{baseline_955_colored.png}
        \label{fig:sub-eleventh}
    \end{subfigure}
    \begin{subfigure}[b]{0.30\linewidth}
        \centering
        \includegraphics[width=\textwidth, height=0.5\linewidth]{original_128_colored_updated.png}
        \label{fig:sub-fourth}
    \end{subfigure}
    \begin{subfigure}[b]{0.30\linewidth}
        \centering
        \includegraphics[width=\textwidth, height=0.5\linewidth]{ours_128_colored_updated.png} 
        \label{fig:sub-sixth}
    \end{subfigure}
    \begin{subfigure}[b]{0.30\linewidth}
        \centering
        \includegraphics[width=\textwidth, height=0.5\linewidth]{baseline_128_colored.png}  
        \label{fig:sub-fifth}
    \end{subfigure}
    \begin{subfigure}[b]{0.30\linewidth}
        \centering
        \includegraphics[width=\textwidth, height=0.5\linewidth]{original_386_colored_updated.png}  
        \label{fig:sub-seventh}
    \end{subfigure}
    \begin{subfigure}[b]{0.30\linewidth}
        \centering
        \includegraphics[width=\textwidth, height=0.5\linewidth]{ours_386_colored_updated.png} 
        \label{fig:sub-ninth}
    \end{subfigure}
    \begin{subfigure}[b]{0.30\linewidth}
        \centering
        \includegraphics[width=\textwidth, height=0.5\linewidth]{baseline_386_colored.png}  
        \label{fig:sub-eighth}
    \end{subfigure}
    \caption{ \textbf{Left:} Input dynamic scan \textbf{Middle}: \dslr{} output. Static structures(red box) \& static occlusions (green box) are retained while reconstructing dynamic occlusion. \textbf{Right:} Best Baseline. Baseline reconstruction fails to discriminate static structures \& static occlusions from dynamic ones. Thus it doesn't retain static structures(red box) and static occlusions(green box).}
    \label{fig:recon}
\end{figure}




\begin{figure*}[htbp]
    \centering
    \includegraphics[width=7in]{Movable-moving-second.pdf}
    \caption{Top view of 3 contiguous interleaving scans (left to right). Segmentation of moving and movable objects using \dslr{}. Movable objects are indicated by pink box, and moving objects by blue. Note the position of the moving object changes while that of the movable remains the same. The black dot in row 2 indicates the origin (0,0). Video demo in Supp.}
    \label{fig:movable-second}
\end{figure*}

\begin{figure}[htbp]
  \centering
\includegraphics[width=1.0\linewidth, height=0.9\linewidth]{ati-movable-comparison.pdf}  
  \caption{Dynamic Object Segmentation comparison on a sparse dataset. \textbf{Top}: Original Dynamic \textbf{Middle}: Dynamic Object segmentation by \dslr{} \textbf{Bottom}: Dynamic Object Segmentation by best baseline. Black dot in rows 2, 3 indicates origin (0,0). For a video demo, please refer Supp.}
\label{fig:ati-moving-movable}
\end{figure}
\subsection{Experimental Setup}
All SLAM experiments were run on an Intel Core i7 CPU. All models were trained using a Nvidia RTX-3090Ti GPU. \dslr{} was trained for 1500 epochs, using a batch size of 64 using RMS optimizer. We use a DCGAN-based generator and a MLP-based discriminator. The learning rate is initialized at 1e-4 for the network. \dslruda{}-based GAN network is trained for 600 epochs using similar parameters as above, after being initialized with CARLA-64 pre-trained weights.
\section{Results}
 

\subsection{
Static Translation using \dslr{}}
\label{sec:chamfer-nu}
CARLA-64 and ARD-16 have correspondence pairs available.
Hence we report the average Chamfer's Distance and Earth Mover's Distance for CARLA-64 and ARD-16 industrial datasets across 5 baselines. (Table \ref{tab:carlaresults}). \dslr{}
performs better at both metrics against all the baselines.
% We also show ablation studies w.r.t. the components of our models over the best-performing  baseline in Table \ref{tab:ablation}. 
The key advantage \dslr{} has is that it's able to overcome the challenges as discussed in (Section \ref{introduction}). \dslr{} is able to \textit{differentiate between static and dynamic(moving and movable) structures}, as well as \textit{between static and dynamic occlusions}. It only omits dynamic structures/occlusions, while preserving static ones. This effect is due to the explicit attention generated on dynamic objects by the couple discriminator and our  strategy of contrasting static LiDAR scan pairs from corresponding dynamic scans (Figure \ref{fig:recon}) using carefully chosen hard negatives. This is also demonstrated in Fig.\ref{fig:movable-first},\ref{fig:movable-second},\ref{fig:ati-moving-movable}, where \dslr{}, successfully segments out moving and movable objects.\\ Although the difference for CARLA and ARD w.r.t the best baseline seems marginal on certain metrics, it is significant enough to generate finer scene details w.r.t moving and movable objects (Fig. \ref{fig:ati-moving-movable}) (Video Demo in Supp.). \dslr{} overcomes the shortcoming of existing baselines and performs better while reconstructing \textit{only} dynamic occlusion holes and retaining all other static structures (Figure \ref{fig:recon}) (Demo in Supp). These results are vital for SLAM in challenging scenarios like sparse datasets on industrial floors (ARD-16).


\setlength{\tabcolsep}{1.5pt}
\begin{table}[htbp]
\centering
\begin{tabular}{|cccccc|}
\hline
\multicolumn{1}{|c|}{\textbf{S-INV}}&\multicolumn{1}{c|}{\textbf{C-AE}}&\multicolumn{1}{c|}{\textbf{C-GAN}}& \multicolumn{1}{c|}{\textbf{C-Net}}&\multicolumn{1}{c|}{\textbf{DSLR}}& \textbf{MOVESe} \\ \hline
\multicolumn{6}{|c|}{CARLA-64 (CD/EMD)}                                 \\ \hline
\multicolumn{1}{|c|}{10.7/947} & \multicolumn{1}{c|}{3.5/311} & \multicolumn{1}{c|}{3.1/281} & \multicolumn{1}{c|}{2.7/553} & \multicolumn{1}{c|}{2.4/226} & \textbf{1.6/203} \\ \hline
\multicolumn{6}{|c|}{ARD-16 (CD/EMD)}                 \\ \hline
\multicolumn{1}{|c|}{7.49/698}     & \multicolumn{1}{c|}{0.95/201}     & \multicolumn{1}{c|}{0.65/195}     & \multicolumn{1}{c|}{3.98/502}     & \multicolumn{1}{c|}{0.37/194}     & \textbf{0.36/128}  \\ \hline
\multicolumn{6}{|c|}{KITTI (LQI)}        \\ \hline
\multicolumn{1}{|c|}{4.8}        & \multicolumn{1}{c|}{1.7}         & \multicolumn{1}{c|}{1.1}         & \multicolumn{1}{c|}{-}            & \multicolumn{1}{c|}{1.1}          & \textbf{0.5}       \\ \hline
\end{tabular}
\caption{Results of our model using Chamfer Distance and EMD for CARLA-64, ARD, and LQI for KITTI with respect to the baselines. Lower is better.}
\label{tab:carlaresults}
\end{table}



\setlength{\tabcolsep}{1.5pt}
\begin{table}[htbp]
\centering
\begin{tabular}{|ccccccc|}
\hline
\multicolumn{1}{|c|}{}           & \multicolumn{3}{c|}{\textbf{Best Baseline}}                                                                                            & \multicolumn{3}{c|}{\dslr{}}                                                                              \\ \hline
\multicolumn{1}{|c|}{Run}        & \multicolumn{1}{c|}{ATE}         & \multicolumn{2}{c|}{RPE}                                                          & \multicolumn{1}{c|}{ATE}            & \multicolumn{2}{c|}{RPE}                                          \\ \hline
\multicolumn{1}{|c|}{}           & \multicolumn{1}{c|}{(m)}         & \multicolumn{1}{c|}{Trans (m)}        & \multicolumn{1}{c|}{Rot ($\theta$)}         & \multicolumn{1}{c|}{(m)}            & \multicolumn{1}{c|}{Trans (m)}      & Rot ($\theta$)                \\ \hline
\multicolumn{1}{|c|}{}           & \multicolumn{6}{c|}{CARLA-64}                                                        \\ \hline
\multicolumn{1}{|c|}{0}          & \multicolumn{1}{c|}{5.67}        & \multicolumn{1}{c|}{0.102}            & \multicolumn{1}{c|}{0.413}                & \multicolumn{1}{c|}{\textbf{1.74}}  & \multicolumn{1}{c|}{\textbf{0.075}} & \textbf{0.412}              \\ \hline
\multicolumn{1}{|c|}{1}          & \multicolumn{1}{c|}{2.71}        & \multicolumn{1}{c|}{\textbf{0.050}}   & \multicolumn{1}{c|}{\textbf{0.390}}       & \multicolumn{1}{c|}{\textbf{1.10}}  & \multicolumn{1}{c|}{0.060}          & \textbf{0.390}              \\ \hline
\multicolumn{1}{|c|}{2}          & \multicolumn{1}{c|}{1.34}        & \multicolumn{1}{c|}{\textbf{0.077}}   & \multicolumn{1}{c|}{\textbf{0.526}}       & \multicolumn{1}{c|}{\textbf{1.11}}  & \multicolumn{1}{c|}{0.078}          & \textbf{0.520}              \\ \hline
\multicolumn{1}{|c|}{3}          & \multicolumn{1}{c|}{6.80}        & \multicolumn{1}{c|}{\textbf{0.136}}   & \multicolumn{1}{c|}{\textbf{0.375}}       & \multicolumn{1}{c|}{\textbf{3.83}}  & \multicolumn{1}{c|}{0.212}          & 0.379                       \\ \hline
\multicolumn{1}{|c|}{}           & \multicolumn{6}{c|}{ARD-16}                      \\ \hline
\multicolumn{1}{|c|}{0}          & \multicolumn{1}{c|}{2.06}        & \multicolumn{1}{c|}{\textbf{0.067}}   & \multicolumn{1}{c|}{1.550}                & \multicolumn{1}{c|}{\textbf{1.97}}  & \multicolumn{1}{c|}{\textbf{0.067}} & \textbf{1.520}              \\ \hline
\multicolumn{6}{|c|}{KITTI}           &                                            \\ \hline
\multicolumn{1}{|c|}{0}          & \multicolumn{1}{c|}{5.22}        & \multicolumn{1}{c|}{0.025}            & \multicolumn{1}{c|}{0.005}                & \multicolumn{1}{c|}{\textbf{2.64}}  & \multicolumn{1}{c|}{\textbf{0.015}} & \textbf{0.003}              \\ \hline
\multicolumn{1}{|c|}{1}          & \multicolumn{1}{c|}{24.64}       & \multicolumn{1}{c|}{0.140}            & \multicolumn{1}{c|}{0.007}                & \multicolumn{1}{c|}{\textbf{4.91}}  & \multicolumn{1}{c|}{\textbf{0.080}} & \textbf{0.003}              \\ \hline
\multicolumn{1}{|c|}{2}          & \multicolumn{1}{c|}{8.98}        & \multicolumn{1}{c|}{0.032}            & \multicolumn{1}{c|}{0.003}                & \multicolumn{1}{c|}{\textbf{5.84}}  & \multicolumn{1}{c|}{\textbf{0.026}} & \textbf{0.002}              \\ \hline
\multicolumn{1}{|c|}{4}          & \multicolumn{1}{c|}{1.85}        & \multicolumn{1}{c|}{0.158}            & \multicolumn{1}{c|}{0.132}                & \multicolumn{1}{c|}{\textbf{0.77}}  & \multicolumn{1}{c|}{\textbf{0.108}} & \textbf{0.124}              \\ \hline
\multicolumn{1}{|c|}{5}          & \multicolumn{1}{c|}{2.05}        & \multicolumn{1}{c|}{0.026}            & \multicolumn{1}{c|}{\textbf{0.004}}       & \multicolumn{1}{c|}{\textbf{1.30}}  & \multicolumn{1}{c|}{\textbf{0.015}} & 0.005                       \\ \hline
\multicolumn{1}{|c|}{6}          & \multicolumn{1}{c|}{4.46}        & \multicolumn{1}{c|}{0.263}            & \multicolumn{1}{c|}{0.025}       & \multicolumn{1}{c|}{\textbf{1.82}}  & \multicolumn{1}{c|}{\textbf{0.060}} & \textbf{0.013}                       \\ \hline
\multicolumn{1}{|c|}{7}          & \multicolumn{1}{c|}{5.23}        & \multicolumn{1}{c|}{0.026}            & \multicolumn{1}{c|}{0.006}                & \multicolumn{1}{c|}{\textbf{2.64}}  & \multicolumn{1}{c|}{\textbf{0.016}} & \textbf{0.003}              \\ \hline
\multicolumn{1}{|c|}{8}          & \multicolumn{1}{c|}{6.70}        & \multicolumn{1}{c|}{0.031}            & \multicolumn{1}{c|}{\textbf{0.006}}       & \multicolumn{1}{c|}{\textbf{2.90}}  & \multicolumn{1}{c|}{\textbf{0.018}} & \textbf{0.006}              \\ \hline
\multicolumn{1}{|c|}{9}          & \multicolumn{1}{c|}{27.96}        & \multicolumn{1}{c|}{0.350}            & \multicolumn{1}{c|}{0.114}                & \multicolumn{1}{c|}{\textbf{4.50}}  & \multicolumn{1}{c|}{\textbf{0.055}} & \textbf{0.043}              \\ \hline
\multicolumn{1}{|c|}{10}         & \multicolumn{1}{c|}{2.45}        & \multicolumn{1}{c|}{0.034}            & \multicolumn{1}{c|}{0.010}                & \multicolumn{1}{c|}{\textbf{0.70}}  & \multicolumn{1}{c|}{\textbf{0.024}} & \textbf{0.008}              \\ \hline
\end{tabular}
\caption{Comparison of SLAM Results on 4 CARLA-64, 1-ARD and 10 KITTI Sequences(00-10). Lower is better.}
\label{tab:slamresults}
\end{table}




\setlength{\tabcolsep}{1.5pt}
\begin{table}[htbp]
\centering
\begin{tabular}{|c|c|c|c|}
\hline
 (Best baseline) & (CoD) & CoD+CL \\ \hline
 \multicolumn{3}{|c|}{{CD/EMD}}         \\ \hline
 2.44/226 & 1.65/206 & \textbf{1.56/203}\\ \hline
\end{tabular}
\caption{Ablation studies show the effect of the couple discriminator(CoD) \& contrastive loss(CL) module in our model when compared with best baseline (DSLR).}
\label{tab:ablation}
\end{table}


For KITTI, we use LQI as a metric due to the unavailability of ground truth static scans. Our method improves to a good extent on the available baselines. This is also evident from the downstream SLAM performance in Table \ref{tab:slamresults}. For all KITTI odometry sequences from 00-10, \dslr{} performs superior to the best-performing baseline, due to the availability of better reconstructed static structures. 
\subsection{Ablation Studies}
\dslr{} has 2 modules: Couple Discriminator(CoD) and Contrastive Loss(CL). We perform ablation studies to show the effect of these modules on \dslr{}. CoD is our pair-discriminator based Relativistic Average Least Squares GAN (col. 2 of Table \ref{tab:ablation}). It is able to bridge a lot of gap between the best baseline and our model. Further, CL along with CoD (Col. 3) gives a good reduction on Chamfer's Distance.
\subsection{Application of \dslr{} for SLAM in dynamic environments}

\label{sec:slamindyenv}
In this section, we demonstrate the effect of \dslr{} on downstream SLAM performance.
We use Google Cartographer \cite{hess2016real}, a LiDAR-based SLAM algorithm to test our proposed approach as a preprocessing step on LiDAR scans before it's consumed by the SLAM algorithm. We compare our results with the best baseline on ARD-16, KITTI and CARLA sequences in Table 
\ref{tab:slamresults}.\\
Although our method doesn't use any segmentation information, we also compare \dslr{} against some strong baselines that use segmentation information in Table \ref{table:slam-seg-results} to show the \textit{benefit of using reconstructed static structures against segmentation-based dynamic point removal.} 
We compare our method against 3 strong baselines (1) original unprocessed dynamic LiDAR scans (Pure-Dynamic) (2) preprocessing to remove dynamic objects using segmentation information 
(3) using the best baseline reconstructions and segmentation information to replace dynamic objects with static reconstructions. \cite{kumar2021dslr}
For (2) and (3), 2 variants(GTS-using ground truth segmentation information, MS- using segmentation information from a trained U-Net) are used.

\begin{table}[htbp]
\centering
\begin{tabular}{|cccc|}
\hline
\multicolumn{1}{|c|}{\textbf{Seg Baselines}}                  & \multicolumn{1}{c|}{\textbf{ATE}}    & \multicolumn{2}{c|}{\textbf{RPE}}              \\ \hline
\multicolumn{1}{|c|}{}                       & \multicolumn{1}{c|}{(m)}    & \multicolumn{1}{c|}{Trans(m)} & Rot   \\ \hline
\multicolumn{4}{|c|}{\textbf{KITTI} }                                         \\ \hline
\multicolumn{1}{|c|}{Pure-Dynamic}           & \multicolumn{1}{c|}{11.809} & \multicolumn{1}{c|}{1.620}    & 1.290 \\ \hline
\multicolumn{1}{|c|}{Detect \& Delete (MS)}  & \multicolumn{1}{c|}{12.490} & \multicolumn{1}{c|}{1.623}    & 1.290 \\ \hline
\multicolumn{1}{|c|}{Detect \& Delete (GTS)} & \multicolumn{1}{c|}{11.458} & \multicolumn{1}{c|}{1.630}    & 1.336 \\ \hline
\multicolumn{1}{|c|}{DSLR-Seg (MS)}          & \multicolumn{1}{c|}{99.642} & \multicolumn{1}{c|}{1.610}    & 1.290 \\ \hline
\multicolumn{1}{|c|}{DSLR-Seg (GTS)}         & \multicolumn{1}{c|}{11.699} & \multicolumn{1}{c|}{1.620}    & 1.295 \\ \hline
\multicolumn{1}{|c|}{MOVESe}                  & \multicolumn{1}{c|}{\textbf{2.80}} & \multicolumn{1}{c|}{\textbf{0.042}}    & \textbf{0.02} \\ \hline
\multicolumn{4}{|c|}{\textbf{CARLA-64 Dataset}}                                                                             \\ \hline
\multicolumn{1}{|c|}{Pure-Dynamic}           & \multicolumn{1}{c|}{10.360} & \multicolumn{1}{c|}{0.056}    & 0.403 \\ \hline
\multicolumn{1}{|c|}{Detect \& Delete (MS)}  & \multicolumn{1}{c|}{10.430} & \multicolumn{1}{c|}{0.060}    & 0.430 \\ \hline
\multicolumn{1}{|c|}{Detect \& Delete (GTS)} & \multicolumn{1}{c|}{11.650} & \multicolumn{1}{c|}{0.063}    & 0.401 \\ \hline
\multicolumn{1}{|c|}{DSLR-Seg (MS)}          & \multicolumn{1}{c|}{10.760} & \multicolumn{1}{c|}{\textbf{0.050}}    & \textbf{0.390} \\ \hline
\multicolumn{1}{|c|}{DSLR-Seg (GTS)}         & \multicolumn{1}{c|}{7.330}  & \multicolumn{1}{c|}{\textbf{0.050}}    & 0.340 \\ \hline
\multicolumn{1}{|c|}{MOVESe}                  & \multicolumn{1}{c|}{\textbf{3.767}}  & \multicolumn{1}{c|}{0.062}    & 0.428 \\ \hline
\multicolumn{4}{|c|}{\textbf{ARD-16 Dataset}}             \\ \hline
\multicolumn{1}{|c|}{Pure-Dynamic}           & \multicolumn{1}{c|}{1.988}  & \multicolumn{1}{c|}{0.068}    & 1.549 \\ \hline
\multicolumn{1}{|c|}{MOVESe}                  & \multicolumn{1}{c|}{\textbf{1.972}}  & \multicolumn{1}{c|}{\textbf{0.067}}    & \textbf{1.525} \\ \hline
\end{tabular}
\caption{Comparison of SLAM  performance of our models against Pure Dynamic as well as segmentation-based baselines.}
\label{table:slam-seg-results}
\end{table}

\begin{figure}[htbp]
  \centering
\includegraphics[width=1.0\linewidth, height=0.5\linewidth]{trajectory_00_01_compare.jpg}  
  \caption{Trajectory comparison between our model and the best baseline for 2 KITTI runs \textbf{Left}: Complex Urban Run \textbf{Right}: Highway Run}
\label{fig:trajectory}
\end{figure}



\subsubsection{Evaluation Metrics}
The trajectory generated using the reconstructed static LiDAR scans accompanied by Cartographer is compared against the ground truth pose estimates using 2 widely used metrics
: Absolute Trajectory Error (ATE) \cite{sturm2012benchmark} and Relative Pose Error (RPE) \cite{sturm2012benchmark}.
\subsubsection{SLAM Analysis}
We compare \dslr{}'s SLAM performance against the best-performing model based on  static translation, with and without segmentation labels.


We would like to point out that despite Google cartographer being efficient at loop closure, we observe high translation errors for the best baseline.
\dslr{} is able to narrow down these translational errors across all the datasets.\\
We find that for runs with loops(complex urban runs) as well as runs without loops (highway), \dslr{} is better than the best baseline as demonstrated in Figure \ref{fig:trajectory}. The key advantage over the baseline here is: As the run progress from the starting point (indicated by a black dot), the baseline (indicated by the green line) starts to drift away from the original trajectory (blue line) faster as compared to ours (red line). This shows that our model is more robust to translation error which denotes that the quality of the occluded static reconstructions is better. Although the difference in translational error might look trivial for some runs in Table \ref{tab:slamresults}, the improvements by \dslr{}, can play a pivotal role in reducing unwanted accidents in restricted environments like industry floors.\\
Another important advantage is shown in Table \ref{table:slam-seg-results} where we perform better than segmentation baselines without the use of annotated segmentation information. This helps us to conclude that generated static structures are adept at increasing the navigation accuracy of SLAM algorithms.

\subsection{Segmentation of Dynamic Objects}

\label{sec:segm_mova_movi}
We qualitatively show that our models is adept at segmenting moving as well as movable objects in LiDAR scans. Using static translation results, we subtract the generated static scans from the original dynamic ones. Given our generated static scans consists of permanent stationary structures only, the subtraction gives us moving as well as movable objects. \dslr{} is capable of segmenting out movable as well as moving objects as shown in Fig \ref{fig:movable-first} and  \ref{fig:movable-second}. The key advantage of \dslr{} is that it is better at accurately identifying only dynamic obstacles and reconstructing the background (unlike existing baselines, which find it difficult to discriminate between static and dynamic obstacles and thus remove static structures as well) as shown in Fig \ref{fig:recon}. 
Moreover, sparse datasets like ARD-16, have very few points falling on dynamic objects. We show that the best baseline model fails to identify the dynamic objects for ARD-16(row 3 of Fig. \ref{fig:ati-moving-movable}). However, our model is successful in segmenting out dynamic objects in such sparse datasets too, as shown in Figure \ref{fig:ati-moving-movable}.
\section{Conclusion}

The paper presents a novel GAN-based adversarial model that helps to segment out moving as well as movable objects using static reconstructions. These reconstructions are accurate and maintain a precise distinction between static and dynamic occlusions and objects.  Moreover, the static structures generated by the model help to improve navigation performance. Our experiments show that SLAM benefits more from by using generated static structures behind occlusion than from removing dynamic points. We achieve all of this in segmentation label-devoid settings.


{\small
\bibliographystyle{unsrt}
\bibliography{egpaper_final}

}

\end{document}